%versi 3 (22-07-2020)

\chapter{Landasan Teori}
\label{chap:teori}

\section{Permasalahan N-Queens}
\label{sec:n-queens}

\subsection{Definisi dan Representasi}
\label{subsec:n-queens-definisi}
Permasalahan N-Queens merupakan salah satu permasalahan klasik dalam bidang ilmu komputer dan matematika diskrit, dengan sejarah panjang yang didokumentasikan secara komprehensif dalam survei yang dilakukan oleh Bell dan Stevens \cite{bell2009survey}. Permasalahan ini pertama kali diperkenalkan oleh Max Bezzel pada tahun 1848 dalam majalah catur Jerman \textit{Berliner Schachzeitung}. Formulasi awal yang diajukan Bezzel hanya membahas kasus papan berukuran $8 \times 8$, dan baru pada tahun 1850 Franz Nauck mempublikasikan seluruh 92 solusi untuk kasus tersebut sekaligus memperumum permasalahan ke papan berukuran $n \times n$.

Pada formulasi umumnya, diberikan sebuah papan catur berukuran $n \times n$ dan tugasnya adalah menempatkan $n$ buah bidak menteri sedemikian rupa sehingga tidak ada dua menteri yang saling menyerang. Sebagaimana dalam aturan permainan catur, sebuah menteri dapat bergerak secara bebas dalam arah horizontal, vertikal, ataupun diagonal dengan jarak tak terbatas, sehingga setiap penempatan menteri harus mempertimbangkan seluruh arah serangan tersebut seperti yang tertera pada Gambar \ref{fig:queenattackpattern}.

\begin{figure}[H]
	\centering
	\includegraphics[width=0.9\linewidth]{queen_attack_pattern.png}
	\caption{Contoh aturan penyerangan bidak menteri pada permainan catur. Silang merah menandakan kotak yang terserang oleh bidak menteri.}
	\label{fig:queenattackpattern}
\end{figure}

Dalam praktiknya, papan dapat direpresentasikan menggunakan beberapa cara. Representasi paling intuitif adalah matriks biner berukuran $n \times n$, di mana nilai $1$ menandakan keberadaan menteri dan nilai $0$ menandakan kotak kosong. Namun, representasi ini tidak efisien dari sudut pandang komputasi karena memungkinkan konfigurasi yang secara trivial tidak valid, seperti dua menteri pada baris yang sama. Oleh sebab itu, representasi yang lebih umum digunakan dalam literatur adalah representasi berbasis permutasi, yaitu sebuah array satu dimensi $\pi = (\pi_1, \pi_2, \dots, \pi_n)$ dengan $\pi_i$ menyatakan kolom tempat menteri pada baris ke-$i$ ditempatkan \cite{glock2022n}. Karena $\pi$ merupakan permutasi dari himpunan $\{1, 2, \dots, n\}$, representasi ini secara implisit menjamin tidak adanya konflik baris maupun konflik kolom, sehingga ruang pencarian yang perlu dieksplorasi berkurang dari $n^n$ menjadi $n!$ konfigurasi.

Dengan demikian, permasalahan N-Queens pada dasarnya dapat dirumuskan sebagai pencarian permutasi $\pi$ yang memenuhi kondisi bahwa tidak ada dua menteri yang berbagi diagonal yang sama. Secara formal, $\pi$ merupakan solusi valid jika dan hanya jika untuk setiap pasangan indeks $i \neq j$ berlaku $|i - j| \neq |\pi_i - \pi_j|$ \cite{glock2022n}. Formulasi ringkas ini menjadi titik berangkat bagi pemodelan permasalahan dalam kerangka kerja yang lebih formal, yang akan dibahas pada Bagian \ref{sec:csp}.

\subsection{Aturan Konflik Queen}
\label{subsec:n-queens-aturan}
Aturan pergerakan bidak menteri pada permainan catur memunculkan tiga jenis konflik yang harus dihindari secara simultan dalam permasalahan N-Queens. Ketiga jenis konflik tersebut bersesuaian dengan tiga arah pergerakan menteri, yaitu horizontal, vertikal, dan diagonal \cite{bell2009survey}.

Konflik horizontal terjadi apabila terdapat dua menteri yang ditempatkan pada baris yang sama. Mengingat sebuah menteri dapat bergerak bebas sepanjang barisnya, kedua menteri tersebut akan saling menyerang tanpa memandang jarak di antara keduanya. Secara analog, konflik vertikal terjadi apabila dua menteri menempati kolom yang sama. Kedua jenis konflik ini bersifat langsung dan mudah diidentifikasi karena hanya melibatkan kesetaraan indeks baris atau kolom.

Konflik diagonal memerlukan analisis yang sedikit lebih cermat. Sebuah menteri dapat menyerang sepanjang dua jenis diagonal, yaitu diagonal utama yang membentang dari kiri atas ke kanan bawah, serta diagonal anti yang membentang dari kanan atas ke kiri bawah. Untuk dua menteri yang ditempatkan pada posisi $(i, \pi_i)$ dan $(j, \pi_j)$, keduanya berada pada diagonal yang sama apabila selisih baris dan selisih kolomnya memiliki nilai absolut yang setara, yaitu $|i - j| = |\pi_i - \pi_j|$. Karakterisasi aritmatika ini memanfaatkan sifat geometris papan: seluruh kotak pada diagonal utama memiliki nilai $i - \pi_i$ yang konstan, sedangkan seluruh kotak pada diagonal anti memiliki nilai $i + \pi_i$ yang konstan \cite{russell2010artificial}.

Berdasarkan ketiga jenis konflik di atas, permasalahan N-Queens dapat dirumuskan secara formal sebagai pencarian penempatan $n$ menteri pada papan $n \times n$ yang memenuhi tiga kendala berikut secara bersamaan: (1) tidak ada dua menteri yang berbagi baris yang sama; (2) tidak ada dua menteri yang berbagi kolom yang sama; dan (3) tidak ada dua menteri yang berbagi diagonal yang sama \cite{russell2010artificial}. Sebagaimana telah dibahas pada Bagian \ref{subsec:n-queens-definisi}, penggunaan representasi permutasi $\pi = (\pi_1, \pi_2, \dots, \pi_n)$ secara implisit menghilangkan dua kendala pertama, sehingga hanya kendala diagonal yang perlu diperiksa secara eksplisit selama proses pencarian solusi.

\begin{figure}[H]
	\centering
	\includegraphics[width=0.35\linewidth]{N_queens_correct.png}
	\caption{Contoh solusi valid permasalahan 8-Queens. Panah hijau menandakan jangkauan serangan salah satu menteri; tidak ada menteri lain yang berada pada baris, kolom, ataupun diagonal yang sama.}
	\label{fig:nqueenscorrect}
\end{figure}

Interaksi antarkendala inilah yang membuat permasalahan N-Queens menarik dari sudut pandang komputasi. Pelanggaran terhadap salah satu kendala mengakibatkan konfigurasi menjadi tidak valid, dan kendala-kendala tersebut tidak dapat diperiksa secara terpisah karena keputusan penempatan satu menteri secara langsung memengaruhi domain posisi valid bagi menteri-menteri lainnya. Sifat saling-bergantung antarkendala ini menjadi dasar bagi pemodelan permasalahan menggunakan kerangka kerja \textit{Constraint Satisfaction Problem} yang akan dibahas pada Bagian \ref{sec:csp}.

\subsection{Kompleksitas dan Relevansi}
\label{subsec:n-queens-kompleksitas}
Dari perspektif komputasi, permasalahan N-Queens menarik karena memiliki struktur yang sederhana namun ruang solusinya sangat besar dan tumbuh secara eksponensial. Apabila tidak ada asumsi tambahan yang digunakan, setiap menteri memiliki $n^2$ kemungkinan posisi pada papan, menghasilkan ruang pencarian awal sebesar $\binom{n^2}{n}$. Penggunaan representasi permutasi sebagaimana dijelaskan pada Bagian \ref{subsec:n-queens-definisi} mempersempit ruang ini menjadi $n!$, namun pertumbuhannya tetap superpolynomial. Sebagai gambaran, untuk $n = 14$ terdapat $14! \approx 8{,}7 \times 10^{10}$ kemungkinan permutasi yang harus dievaluasi, jumlah yang jauh melampaui kemampuan komputasi \textit{brute-force} konvensional pada perangkat keras umum.

Meskipun ruang pencarian tumbuh secara eksponensial, perlu ditekankan bahwa permasalahan N-Queens dalam formulasi dasarnya bukanlah permasalahan yang sulit secara teoretis. Hoffman, Loessi, dan Moore \cite{hoffman1969constructions} sebagaimana dirangkum dalam survei Bell dan Stevens \cite{bell2009survey} menunjukkan bahwa terdapat konstruksi aritmatika eksplisit yang dapat menghasilkan satu solusi N-Queens untuk setiap $n \geq 4$ dalam waktu linier $O(n)$ tanpa perlu melakukan pencarian sama sekali. Dengan demikian, persoalan keputusan ``apakah solusi N-Queens ada untuk papan $n \times n$'' dapat dijawab dalam waktu konstan: solusinya selalu ada kecuali untuk $n = 2$ dan $n = 3$.

Karakteristik kompleksitas N-Queens telah menjadi topik diskusi yang panjang dalam komunitas \textit{Artificial Intelligence}, khususnya terkait kelayakannya sebagai \textit{benchmark} algoritma pencarian. Gent, Jefferson, dan Nightingale \cite{gent2017complexity} memberikan analisis komprehensif terhadap persoalan ini dan menunjukkan bahwa permasalahan N-Queens dalam bentuk dasarnya tidak ideal sebagai tolok ukur kesulitan algoritma, mengingat solusinya dapat dikonstruksi secara langsung tanpa proses pencarian. Sebaliknya, mereka membuktikan bahwa varian \textit{N-Queens Completion}, yaitu permasalahan menyelesaikan konfigurasi parsial yang telah diberikan sebelumnya, tergolong NP-Complete dan \#P-Complete. Hasil ini mengindikasikan bahwa varian-varian N-Queens dengan kendala tambahan atau struktur papan ireguler memiliki kompleksitas teoritis yang setara dengan permasalahan-permasalahan sulit klasik dalam ilmu komputer.

Permasalahan \textit{Colored Queens} yang menjadi fokus tugas akhir ini memiliki struktur kendala yang jauh lebih kompleks dibandingkan N-Queens klasik, termasuk pembagian sektor warna dan kendala \textit{adjacency} yang tidak dapat diselesaikan dengan konstruksi aritmatika sederhana. Dengan kata lain, meskipun teknik konstruktif tersedia untuk N-Queens dasar, teknik tersebut tidak dapat diperumum ke varian-varian dengan kendala ireguler. Hal inilah yang menjadikan pendekatan algoritmik berbasis pencarian, baik yang bersifat sistematis maupun \textit{metaheuristic}, tetap relevan dan menjadi fokus pembahasan pada bab-bab selanjutnya.

\section{Colored Queens Problem}
\label{sec:colored-queens}

\subsection{Definisi dan Aturan Permainan}
\label{subsec:colored-queens-definisi}
Permasalahan \textit{Colored Queens} merupakan sebuah \textit{puzzle} logika berbasis papan yang dipopulerkan melalui permainan \textit{Queens} pada platform LinkedIn \cite{linkedin2024queens}. Secara struktural, permasalahan ini memiliki kesamaan yang kuat dengan jenis \textit{puzzle} penempatan objek yang dikenal sebagai \textit{Star Battle}, yang pertama kali dirancang oleh Hans Eendebak untuk \textit{World Puzzle Championship} tahun 2003 \cite{gmpuzzles2013starbattle}. Dalam varian yang dibahas pada tugas akhir ini, permainan menggunakan papan berukuran $n \times n$ yang dibagi menjadi $n$ buah sektor berwarna (\textit{colored regions}) dengan bentuk yang umumnya tidak beraturan (\textit{irregular}). Setiap sektor warna membentuk satu wilayah yang terhubung (\textit{connected region}) pada papan, artinya seluruh sel dalam satu sektor dapat dijangkau dari sel lainnya melalui tetangga horizontal atau vertikal tanpa melewati sel dari sektor lain. Contoh papan permainan beserta solusinya diperlihatkan pada Gambar \ref{fig:coloredqueens}.

\begin{figure}[H]
	\centering
	\begin{subfigure}[b]{0.45\linewidth}
		\centering
		\includegraphics[width=\linewidth]{Queens_unsolved.png}
		\caption{Papan awal}
		\label{fig:cq-unsolved}
	\end{subfigure}
	\hfill
	\begin{subfigure}[b]{0.45\linewidth}
		\centering
		\includegraphics[width=\linewidth]{Queens_solved.png}
		\caption{Papan yang telah diselesaikan}
		\label{fig:cq-solved}
	\end{subfigure}
	\caption{Contoh papan permainan \textit{Colored Queens} berukuran $8 \times 8$ dengan 8 sektor warna. Setiap sektor membentuk wilayah terhubung, dan solusi menempatkan tepat satu menteri per baris, kolom, dan sektor warna tanpa ada dua menteri yang bertetangga.}
	\label{fig:coloredqueens}
\end{figure}

Tujuan permainan adalah menempatkan tepat $n$ buah bidak menteri pada papan sedemikian rupa sehingga seluruh aturan berikut terpenuhi secara bersamaan:

\begin{enumerate}
	\item Setiap baris memuat tepat satu menteri.
	\item Setiap kolom memuat tepat satu menteri.
	\item Setiap sektor warna memuat tepat satu menteri.
	\item Tidak ada dua menteri yang menempati sel-sel yang bertetangga (\textit{adjacent}), termasuk tetangga diagonal. Secara formal, dua sel $(r_1, c_1)$ dan $(r_2, c_2)$ dikatakan bertetangga apabila $|r_1 - r_2| \leq 1$ dan $|c_1 - c_2| \leq 1$.
\end{enumerate}

Perlu dicatat bahwa, berbeda dengan permasalahan N-Queens klasik yang dibahas pada Bagian \ref{sec:n-queens}, bidak menteri dalam \textit{Colored Queens} \textbf{tidak} menyerang sepanjang diagonal. Serangan hanya berlaku dalam arah horizontal dan vertikal, sehingga kendala diagonal digantikan oleh kendala \textit{adjacency} yang bersifat lokal (hanya memengaruhi delapan sel di sekeliling menteri). Perbedaan ini mengubah karakteristik ruang pencarian secara fundamental, karena kendala \textit{adjacency} tidak dapat direduksi menjadi perbandingan aritmatika sederhana sebagaimana kendala diagonal pada N-Queens.

Selain itu, keberadaan sektor warna yang berbentuk ireguler menambah dimensi kendala yang tidak dimiliki oleh N-Queens. Pada N-Queens, satu-satunya entitas pembatas adalah baris, kolom, dan diagonal yang semuanya memiliki struktur geometris teratur. Pada \textit{Colored Queens}, sektor warna memiliki bentuk yang bervariasi antarpapan dan tidak mengikuti pola geometris yang dapat diprediksi, sehingga setiap instansi permainan memerlukan analisis kendala yang berbeda.

\subsection{Kompleksitas dan Struktur Masalah}
\label{subsec:colored-queens-kompleksitas}
Permasalahan \textit{Colored Queens} memiliki struktur kendala berlapis (\textit{multi-layered constraints}) yang membedakannya secara fundamental dari N-Queens klasik. Sebagaimana diformulasikan oleh Mata Ali dan Mencia \cite{ali2026quboformulationgeneralizedlinkedin}, permainan ini merupakan varian dari N-Queens yang melonggarkan kendala diagonal namun menambahkan kendala sektor warna dan \textit{adjacency}. Kombinasi kendala-kendala tersebut menghasilkan interaksi yang lebih kompleks dibandingkan N-Queens, karena setiap keputusan penempatan menteri secara simultan memengaruhi tiga domain pembatas yang berbeda: baris/kolom, sektor warna, dan lingkungan tetangga.

Lapisan pertama adalah kendala baris dan kolom, yang secara struktural identik dengan N-Queens dan membatasi setiap baris serta kolom untuk memuat tepat satu menteri. Lapisan kedua adalah kendala sektor warna, yang mengharuskan setiap sektor memiliki tepat satu menteri. Berbeda dengan kendala baris dan kolom yang bersifat reguler dan simetris, sektor warna memiliki bentuk ireguler yang bervariasi antarpapan. Implikasinya, ukuran domain untuk setiap sektor (jumlah sel yang tersedia) tidak seragam, dan beberapa sektor mungkin memiliki domain yang jauh lebih kecil dibandingkan sektor lainnya. Hal ini menciptakan ketidakseimbangan dalam proses pencarian: penempatan menteri pada sektor berukuran kecil memiliki dampak propagasi yang berbeda dibandingkan penempatan pada sektor berukuran besar.

Lapisan ketiga adalah kendala \textit{adjacency}, yang melarang dua menteri menempati sel-sel yang bertetangga. Kendala ini bersifat lokal dan simetris, namun efeknya bersifat global karena penempatan satu menteri secara langsung mengeliminasi hingga delapan sel tambahan dari domain menteri-menteri lainnya. Pada papan berukuran kecil ($n \leq 8$), eliminasi ini sangat signifikan secara proporsional terhadap total sel yang tersedia.

Interaksi antara ketiga lapisan kendala tersebut menciptakan beberapa tantangan komputasi. Pertama, tidak ada representasi permutasi sederhana sebagaimana pada N-Queens, karena variabel terkait erat dengan sektor warna yang bentuknya tidak teratur. Kedua, kendala \textit{adjacency} tidak dapat diekspresikan sebagai ketidaksamaan aritmatika global (seperti $|i - j| \neq |\pi_i - \pi_j|$ pada N-Queens), melainkan harus diperiksa secara lokal untuk setiap pasangan sel. Ketiga, bentuk ireguler sektor warna menyebabkan \textit{constraint graph} yang dihasilkan tidak memiliki struktur simetris yang dapat dieksploitasi oleh teknik konstruktif.

Dari sudut pandang kompleksitas, meskipun belum terdapat pembuktian formal bahwa \textit{Colored Queens} termasuk dalam kelas NP-Complete, struktur permasalahannya memiliki kemiripan dengan \textit{N-Queens Completion} yang telah terbukti NP-Complete \cite{gent2017complexity}. Keduanya melibatkan kendala tambahan yang merusak keteraturan geometris N-Queens dasar, sehingga teknik konstruktif tidak dapat diterapkan. Kompleksitas inilah yang menjadikan permasalahan ini menarik untuk diselesaikan secara algoritmik menggunakan pendekatan berbasis pencarian.

\subsection{Posisi dan Relevansi Penelitian}
\label{subsec:colored-queens-relevansi}
Permasalahan \textit{Colored Queens} merupakan topik yang relatif baru dalam literatur akademis. Meskipun permainan ini telah menarik perhatian luas sejak diperkenalkan oleh LinkedIn sebagai \textit{puzzle} harian \cite{linkedin2024queens}, kajian ilmiah formal terhadap permasalahan ini masih sangat terbatas. Sejauh tinjauan pustaka yang dilakukan, satu-satunya publikasi akademis yang secara langsung memformulasikan permainan LinkedIn Queens adalah karya Mata Ali dan Mencia \cite{ali2026quboformulationgeneralizedlinkedin}, yang mengusulkan formulasi \textit{Quadratic Unconstrained Binary Optimization} (QUBO) untuk menyelesaikan generalisasi permainan tersebut pada perangkat keras kuantum. Pendekatan-pendekatan lain yang tersedia di komunitas pengembang umumnya bersifat informal dan tidak dipublikasikan melalui kanal akademis, seperti implementasi \textit{backtracking} sederhana, formulasi \textit{integer programming}, serta penggunaan \textit{SAT/SMT solver}.

Celah penelitian (\textit{research gap}) yang teridentifikasi dari tinjauan ini adalah sebagai berikut. Pertama, belum terdapat kajian yang memodelkan \textit{Colored Queens} secara eksplisit sebagai \textit{Constraint Satisfaction Problem} (CSP) dan menyelesaikannya menggunakan teknik-teknik klasik CSP seperti \textit{backtracking} dengan \textit{forward checking} dan \textit{arc consistency}. Kedua, belum ditemukan penelitian yang membandingkan pendekatan sistematis (\textit{complete}) dengan pendekatan \textit{metaheuristic} untuk permasalahan ini. Ketiga, analisis empiris mengenai pengaruh ukuran papan terhadap performa algoritma penyelesaian juga belum dilakukan secara formal.

Tugas akhir ini bertujuan untuk mengisi celah tersebut dengan memberikan dua kontribusi utama. Kontribusi pertama adalah pemodelan \textit{Colored Queens} sebagai CSP dan penyelesaiannya menggunakan \textit{backtracking} yang diperkuat dengan \textit{forward checking} dan \textit{Arc Consistency} (AC-3). Kontribusi kedua adalah penerapan \textit{Particle Swarm Optimization} (PSO) yang diadaptasi untuk ruang pencarian diskrit sebagai pendekatan alternatif. Perbandingan antara kedua pendekatan tersebut diharapkan memberikan wawasan mengenai \textit{trade-off} antara kelengkapan solusi dan efisiensi waktu, serta perilaku masing-masing algoritma pada papan berukuran kecil hingga besar.

\section{Constraint Satisfaction Problem}
\label{sec:csp}

\subsection{Definisi dan Komponen CSP}
\label{subsec:csp-definisi}
\textit{Constraint Satisfaction Problem} (CSP) merupakan kerangka kerja formal untuk merepresentasikan dan menyelesaikan permasalahan yang melibatkan sekumpulan variabel dengan batasan-batasan yang harus dipenuhi secara bersamaan. Secara formal, sebuah CSP didefinisikan sebagai tripel $\mathcal{P} = \langle X, D, C \rangle$ \cite{russell2010artificial} di mana:

\begin{itemize}
	\item $X = \{x_1, x_2, \dots, x_n\}$ adalah himpunan variabel yang nilainya harus ditentukan.
	\item $D = \{D_1, D_2, \dots, D_n\}$ adalah himpunan domain, di mana $D_i$ menyatakan himpunan nilai yang dapat diambil oleh variabel $x_i$.
	\item $C = \{C_1, C_2, \dots, C_m\}$ adalah himpunan kendala, di mana setiap kendala $C_j$ mendefinisikan kombinasi nilai yang diperbolehkan untuk suatu subhimpunan variabel tertentu.
\end{itemize}

Sebuah \textit{assignment} adalah pemetaan dari satu atau lebih variabel ke nilai-nilai dalam domainnya. \textit{Assignment} disebut \textit{consistent} apabila tidak melanggar kendala manapun, dan disebut \textit{complete} apabila seluruh variabel telah diberi nilai. Solusi dari sebuah CSP adalah \textit{complete assignment} yang bersifat \textit{consistent}, yaitu seluruh variabel telah diberi nilai dan seluruh kendala terpenuhi \cite{russell2010artificial}.

Kajian formal mengenai CSP berawal dari karya Montanari pada tahun 1974 yang memperkenalkan kerangka jaringan kendala (\textit{constraint network}), dan dilanjutkan oleh Mackworth pada tahun 1977 yang mengembangkan algoritma \textit{arc consistency} sebagai teknik dasar untuk mereduksi ruang pencarian \cite{montanari1974networks}. Sejak saat itu, CSP telah menjadi paradigma fundamental dalam \textit{Artificial Intelligence} dan digunakan secara luas untuk memodelkan berbagai permasalahan seperti penjadwalan, pewarnaan graf, dan penyelesaian \textit{puzzle} \cite{russell2010artificial}.

\subsection{Teknik Representasi}
\label{subsec:csp-representasi}

Sebuah CSP dapat direpresentasikan secara visual menggunakan \textit{constraint graph}, yaitu sebuah graf di mana setiap simpul merepresentasikan variabel dan setiap sisi menghubungkan dua variabel yang terlibat dalam suatu kendala bersama \cite{russell2010artificial}. Representasi ini berguna untuk menganalisis struktur ketergantungan antarvariabel serta menentukan strategi penyelesaian yang tepat.

Kendala dalam CSP dapat diklasifikasikan berdasarkan jumlah variabel yang terlibat. Kendala \textit{unary} melibatkan satu variabel dan membatasi domain variabel tersebut secara langsung, misalnya $x_1 \neq 3$. Kendala \textit{binary} melibatkan dua variabel dan membatasi kombinasi nilai yang diperbolehkan di antara keduanya, misalnya $x_1 \neq x_2$. Dalam praktiknya, sebagian besar CSP dapat ditransformasikan ke dalam bentuk yang hanya mengandung kendala \textit{binary} tanpa kehilangan informasi \cite{russell2010artificial}, dan banyak algoritma penyelesaian CSP dirancang secara khusus untuk menangani kendala jenis ini.

Struktur \textit{constraint graph} memiliki implikasi langsung terhadap kompleksitas penyelesaian. Apabila graf memiliki struktur pohon (\textit{tree-structured}), CSP dapat diselesaikan dalam waktu polinomial. Sebaliknya, graf yang padat (\textit{dense}) dengan banyak sisi mengindikasikan interaksi kendala yang tinggi dan umumnya memerlukan teknik penyelesaian yang lebih canggih \cite{russell2010artificial}.

\subsection{Pemodelan Colored Queens sebagai CSP}
\label{subsec:csp-colored-queens}

Permasalahan \textit{Colored Queens} dapat dimodelkan secara natural sebagai CSP dengan mendefinisikan komponen-komponen $\langle X, D, C \rangle$ sebagai berikut.

Himpunan variabel $X = \{x_1, x_2, \dots, x_n\}$ merepresentasikan $n$ buah sektor warna pada papan, di mana setiap variabel $x_i$ bersesuaian dengan sektor warna ke-$i$. Nilai yang diberikan pada $x_i$ menyatakan posisi sel tempat menteri ditempatkan dalam sektor tersebut.

Domain $D_i$ untuk setiap variabel $x_i$ adalah himpunan seluruh sel yang termasuk dalam sektor warna ke-$i$. Secara formal, apabila sektor warna ke-$i$ terdiri dari sel-sel $\{(r_1, c_1), (r_2, c_2), \dots, (r_k, c_k)\}$, maka $D_i = \{(r_1, c_1), (r_2, c_2), \dots, (r_k, c_k)\}$. Perlu dicatat bahwa ukuran domain bervariasi antarsektor karena bentuk sektor yang ireguler, sehingga $|D_i| \neq |D_j|$ pada umumnya.

Himpunan kendala $C$ terdiri dari empat jenis kendala yang bersesuaian dengan aturan permainan:

\begin{enumerate}
	\item \textbf{Kendala baris}: untuk setiap pasangan variabel $x_i, x_j$ dengan $i \neq j$, apabila $x_i = (r_i, c_i)$ dan $x_j = (r_j, c_j)$, maka $r_i \neq r_j$.
	\item \textbf{Kendala kolom}: untuk setiap pasangan variabel $x_i, x_j$ dengan $i \neq j$, berlaku $c_i \neq c_j$.
	\item \textbf{Kendala sektor warna}: dipenuhi secara implisit oleh definisi variabel, karena setiap variabel hanya dapat mengambil nilai dari sektornya sendiri.
	\item \textbf{Kendala \textit{adjacency}}: untuk setiap pasangan variabel $x_i, x_j$ dengan $i \neq j$, apabila $x_i = (r_i, c_i)$ dan $x_j = (r_j, c_j)$, maka $|r_i - r_j| > 1$ atau $|c_i - c_j| > 1$.
\end{enumerate}

Seluruh kendala di atas bersifat \textit{binary}, yaitu melibatkan tepat dua variabel, sehingga permasalahan \textit{Colored Queens} menghasilkan \textit{constraint graph} yang bersifat \textit{complete}—setiap pasangan variabel terhubung oleh setidaknya satu kendala. Karakteristik ini mengindikasikan tingkat interaksi kendala yang tinggi dan memperkuat argumen bahwa teknik propagasi kendala diperlukan untuk mereduksi ruang pencarian secara efektif sebelum atau selama proses pencarian solusi. Teknik-teknik penyelesaian tersebut akan dibahas pada bagian berikutnya.

\section{Teknik Penyelesaian CSP}
\label{sec:teknik-penyelesaian}

\subsection{Algoritma Sistematis}
\label{subsec:algoritma-sistematis}

\subsubsection{Backtracking Search}
\label{subsubsec:backtracking}

\subsubsection{Forward Checking}
\label{subsubsec:forward-checking}

\subsubsection{Arc Consistency (AC-3)}
\label{subsubsec:ac3}

\subsection{Metaheuristic}
\label{subsec:metaheuristic}

\subsubsection{Konsep Umum Metaheuristic}
\label{subsubsec:metaheuristic-konsep}

\subsubsection{Representasi Solusi Diskrit}
\label{subsubsec:representasi-diskrit}

\subsubsection{Particle Swarm Optimization}
\label{subsubsec:pso}


\section{Perbandingan Pendekatan}
\label{sec:perbandingan}

\subsection{Complete vs Metaheuristic}
\label{subsec:perbandingan-complete-meta}

\subsection{Backtracking+FC+AC-3 vs PSO}
\label{subsec:perbandingan-bt-pso}

\subsection{Justifikasi Pemilihan Metode}
\label{subsec:justifikasi-metode}